\documentclass[12pt]{article}
\usepackage[utf8]{inputenc}
\usepackage[T1]{fontenc}
\usepackage{hyperref}
\title{Stod Bhoti Language Resources for K-12 Curriculum Development}
\author{Compiled by Claude Sonnet 5.5 (claude-sonnet-5-5) for Andr\'as Kornai}
\date{}
\begin{document}
\maketitle

\section{Introduction}
Stod Bhoti (ISO 639-3: sbu; sTodpa, also Lahul Bhoti) is a Tibetic language of Lahaul and Spiti district, Himachal Pradesh, India. Dialects are Stod proper (Kolong), Khoksar and Mayar. It is close to Spiti Bhoti (about 74\% lexical similarity). It should not be confused with the Stod subdialect of Zangskari or with Bhoti Kinnauri (Nyamkat). Some grammatical description exists, mostly by the native-speaker linguist Tobdan. No school use was found.

\section{Primary Educational Resources}
No primer or textbook was found.

\section{Linguistic and Reference Materials}
\begin{itemize}
\item Wikipedia, Stod Bhoti: overview, dialects and references, including Tobdan's \emph{A Grammar of sTodpa} (2015) and Sharma (1992). \url{https://en.wikipedia.org/wiki/Stod_Bhoti}
\item Tobdan (2020), \emph{Stodpa language of Lahul in the Himalaya}, Kaveri Books (Glottolog record). \url{https://glottolog.org/resource/reference/id/589790}
\item Tobdan (1990), \emph{Bhoti Parichaya -- Vayakaran aur Anuvad}, a Hindi grammar and translation (Glottolog record). \url{https://glottolog.org/resource/reference/id/577955}
\item Crossgram typological entry. \url{https://crossgram.clld.org/languages/stod1241}
\item PHOIBLE phonological inventory. \url{https://phoible.org/languages/stod1241}
\item Numeral data. \url{https://lingweb.eva.mpg.de/channumerals/Stod-Bhotii.htm}
\end{itemize}
No dictionary was found.

\section{Religious and Community Resources}
None found.

\section{Digital and Audio Resources}
No audio, video or corpus resources were found.

\section{Curriculum Development Context}
Tobdan's grammars are the best starting points. The language is written in the Tibetan script in related varieties, but no orthography for Stod Bhoti was confirmed. Hindi and English are the school languages. A curriculum would need new materials, and Tobdan's work and the local community would be the first contacts.
\end{document}